\begin{proposition}[A partial cusp filling]\hypertarget{tag-030M}{}\label{push:src:B.4.18}
\label{push:B418}
\StatementRecord{push:src:B.4.18}{proved}{030M}{\SeriesRef{push:src:B.3.1}, \SeriesRef{push:src:B.3.2}, \SeriesRef{push:src:B.4.9}, \SeriesRef{push:src:A.2.1}, \SeriesRef{push:src:A.1.2}, \SeriesRef{push:src:B.3.3}}

Put
\[
 B^\times=(\Delta_q\times\Delta_r)\setminus\{(0,0)\}.
\]
There is a proper Hausdorff analytic family over
$B^\times$ with smooth total fourfold, flat over $r$,
whose restriction to $r\ne0$ is the prescribed
cusp family after $w=r^6$, and whose reduced
normal-crossings $r$-central fibre over $q\ne0$ has
component deck permutation the shift $-1$ on $\Z/12$
for the clockwise cusp generator $g_0$ of \ref{push:src:B.3.1}.
Equivalently, the positive $q$ loop has shift $+1$.
In particular its $q$-fibre for every $r\ne0$ is the
reduced Mumford fibre, while its other axis has
the component transport required by \ref{push:src:B.3.2}.

The overlap is the displayed holomorphic translation with its
affine marking on the two specified base opens.
\end{proposition}



\begin{proof}
\begin{enumerate}[label=\arabic*.,leftmargin=*]
\item \emph{The fractional deck and its exact comparison.}
On the logarithmic cover of
$U=\Delta_q^*\times\Delta_r^*$ write
\[
 s=\frac{\log q}{2\pi i},\qquad
 t=\frac{\log r}{2\pi i},\qquad
 C(q)=\beta^{\rm part}(q)+s.
\]
The functions $\mu(q)$ and $C(q)$ are single-valued
and holomorphic at $q=0$. After the affine gauge of
\ref{push:src:B.4.9} the family is its Jacobian, with period
\begin{equation}
 \boldsymbol\gamma(s,t)
       =(6\mu(q),C(q)-s-12t),\qquad
 \boldsymbol\delta=(0,1).
\label{pub:I:eq:B-4-18-1}
\end{equation}
The other periods are
$\mathbf u=(\tau,\mu)$ and $\mathbf w=(1,0)$.
Under $s\mapsto s+1$, they change by
$\boldsymbol\gamma\mapsto\boldsymbol\gamma-
\boldsymbol\delta$ and $\mathbf u\mapsto\mathbf u+\mathbf w$.
Under $t\mapsto t+1$, the period $\boldsymbol\gamma$
changes by $-12\boldsymbol\delta$.
These are integral changes of the period basis; the
ambient additive $\C^2$ coordinate is unchanged.

\item Set
\[
 b(s,t)=\frac1{12}\boldsymbol\gamma(s,t)
                         -\frac1{24}\boldsymbol\delta.
\]
The two exact identities
\begin{equation}
 \sum_{k=0}^{11}b(s+k,t)
       =\boldsymbol\gamma(s,t)-6\boldsymbol\delta,
 \qquad
 b(s,t+1)-b(s,t)=-\boldsymbol\delta
\label{pub:I:eq:B-4-18-2}
\end{equation}
show that this gives a genuine order-twelve affine
$q$-deck action on the root cover $q=z^{12}$,
well defined around the $r$ loop. The first right-hand
side is a period. The sum of the twelve fractions
$\boldsymbol\gamma(s+k,t)/12$ is
$\boldsymbol\gamma(s,t)-11\boldsymbol\delta/2$.
The twelve constant terms $-\boldsymbol\delta/24$
contribute $-\boldsymbol\delta/2$, giving the integral period
$\boldsymbol\gamma(s,t)-6\boldsymbol\delta$.

\item There is an explicit holomorphic trivialization on $U$.
Put $L=s+12t$ and
\begin{equation}
 F(s,t)=
 \left(\frac{\mu(q)L}{2},\,
       \frac{C(q)L}{12}-\frac{L^2}{24}+6t\right).
\label{pub:I:eq:B-4-18-3}
\end{equation}
Direct subtraction gives
\begin{equation}
 F(s+1,t)-F(s,t)=b(s,t),\qquad
 F(s,t+1)-F(s,t)=\boldsymbol\gamma(s,t).
\label{pub:I:eq:B-4-18-4}
\end{equation}
The second difference is the displayed integral period, so the
translation descends to the torus.
Thus $\zeta\mapsto\zeta+F$ is a holomorphic
isomorphism from the untwisted Jacobian to the
$+b$ affine torus torsor over $U$; its inverse
$\zeta\mapsto\zeta-F$ trivializes that torsor.
Replacing $b,F$ by $-b,-F$ gives the opposite
component rotation. If $\sigma$ denotes the 
affine section of \ref{push:src:B.4.9}, the map from the 
family to the $+b$ torsor is explicitly
$\zeta\mapsto\zeta-\sigma+F$.
This map preserves the affine cusp marking.

\item \emph{A rotating compactification before changing the
$q$-boundary.}
In the root-cover model, the two period shears in cocharacter coordinates
$(y_1,y_2,h_z,h_r)$ are
\[
 \begin{aligned}
 G_1 &: (y_1,y_2)\longmapsto(y_1+12h_z,y_2),\\
 G_2 &: (y_1,y_2)\longmapsto
                         (y_1,y_2+12h_z+12h_r).
 \end{aligned}
\]
The positive fractional deck has the shear
\begin{equation}
 H:(y_1,y_2)\longmapsto(y_1,y_2+h_z+h_r),
 \qquad H^{12}=G_2,
\label{pub:I:eq:B-4-18-5}
\end{equation}
and rotates $z$ by $e^{2\pi i/12}$.
It commutes with $G_1$. Its unit factors are holomorphic
functions of $q$. For example, in the convention
$W=e^{-2\pi i\zeta_2}$, it multiplies $W$ by
\[
 r z\,e^{-2\pi i C(q)/12}e^{\pi i/12}
\]
and the first multiplicative coordinate by
$e^{\pi i\mu(q)}$. The twelve successive root rotations
contribute the factor
$e^{2\pi i(0+\cdots+11)/12}=-1$, and the scalar
correction contributes $e^{\pi i}=-1$.
They cancel. The twelfth iterate is consequently
the second period, including its unit, as
also follows additively from \eqref{pub:I:eq:B-4-18-2}.

\item Use the unit-square cusp triangulation, with twelve
consecutive bottom cells in the $y_1$ direction and
the same bridge choice in every unit prism. Its
translates have bottom cell indices $(i,j)$ and top
interval $[j,j+1]$. Thus $H$ increases the bottom
second index and the top index by one, while $G_1$
shifts the first index by twelve. The fan is invariant
under both, and therefore under $G_2$.
For each bottom cell its three maximal cones are
the determinant-one cones of \ref{push:src:A.2.1}. The usual
compact-core construction of \ref{push:src:A.1.2}, now on the
$I_{12}$ elliptic root model and with twelve
plumbing intervals before the $G_2$ quotient,
gives a proper model on the root cover. Its
unfolded charts are precisely this fan. The compact core and collars therefore prove properness and
proper discontinuity.

\item The affine deck extends to this model because its
unit multiplication preserves every toric chart
and its shear preserves the fan. It is free.
On the $r$-central divisor a point lies on at most
two adjacent component labels in $\Z/12$; no
nonidentity translation can preserve that set.
This covers their intersection over the common
base origin as well. On the $q$-central divisor
with $r\ne0$, the bottom triangle has second
component indices contained in $\{j,j+1\}$,
again not preserved by a nonzero translation
modulo twelve. Off both axes the root-base action
is free. Its finite quotient is therefore smooth,
proper over $(q,r)$, flat over $r$, and has a
reduced normal-crossings $r$-central divisor.

\item Keeping the same positive root-base rotation and using translation
$-b$ gives valuation shear $H^{-1}$ and twelfth iterate
$G_2^{-1}$. This action preserves the same fan and has the same
freeness properties. For the family in the statement we use
translation by $b$ and the clockwise cusp generator
$g_0:s\mapsto s-1$. Thus the
positive fractional deck $b$ has shift $+1$ for
$s\mapsto s+1$ and shift $-1$ for $g_0$, exactly
as required by \ref{push:src:B.3.2}. The finite shifts $4,-3$
and this cusp shift all refer to the same clockwise
deck-identification convention.

\item Denote this rotating model by $\mathcal R$.
Its $q$-fibre for $r\ne0$ has multiplicity twelve.
Indeed the original $z$-central divisor is reduced
and the finite action has trivial generic isotropy
there, while $q=z^{12}$. This is the $q$-multiplicity computed in \ref{push:src:B.3.3}.

\item \emph{The two-open gluing.}
On the base open set $q\ne0$ use $\mathcal R$ with
the positive fractional deck $b$. On the base open set $r\ne0$
use the reduced Mumford cusp family.
Their common base is $U$, where \eqref{pub:I:eq:B-4-18-3}--
\eqref{pub:I:eq:B-4-18-4}, composed with \ref{push:src:B.4.9}, give the specified
holomorphic isomorphism. These two open sets cover
$B^\times$, and there is no additional triple
overlap. Glue by that isomorphism.

\item The result is Hausdorff: points over one base
point are separated in one of the two given
Hausdorff models, and points with different
base images are separated in the base.
Properness is local on the target, so follows
from the properness of those two restrictions.
Smoothness and flatness over $r$ are also local.
The restriction over $r\ne0$ is, by construction,
the prescribed cusp including its reduced
$q$-fibre. Over $r=0,q\ne0$, the construction has
the required twelve-cycle and rotation $-1$.
This proves every asserted property.

\end{enumerate}
\end{proof}